\makeatletter
\def\rap@fontpath{../../../Common/}
\makeatother
\documentclass{../../../Common/Latex/Templates/rapport}
% ---------------------------------------------------------------------------
% Shared settings for the H2026 General Assembly document set.
%
% Usage, directly after \documentclass{../../../Common/Latex/Templates/rapport}
% in a document one folder below this file:
%   % ---------------------------------------------------------------------------
% Shared settings for the H2026 General Assembly document set.
%
% Usage, directly after \documentclass{../../../Common/Latex/Templates/rapport}
% in a document one folder below this file:
%   % ---------------------------------------------------------------------------
% Shared settings for the H2026 General Assembly document set.
%
% Usage, directly after \documentclass{../../../Common/Latex/Templates/rapport}
% in a document one folder below this file:
%   \input{../ga-common}
%
% Image paths are relative to the compiled document's folder, so every
% document using this file must sit one folder below it.
% ---------------------------------------------------------------------------

% Meeting details, used on covers and in each document's ID block.
\newcommand{\meetingdate}{30 September 2026}
\newcommand{\meetingtime}{12:00}
\newcommand{\meetinglocation}{Festsalen, NMBU}
\newcommand{\meetingline}{\textbf{Meeting:} \meetingdate, \meetingtime, \meetinglocation}

% Header: moon logo on the left, full-text logo on the right. \projectname is
% still set because the cover title uses it.
\projectname{Eclipse NMBU}
\projectlogo{../../../Common/Eclipse Images/EclipseMoon}
\orglogo{../../../Common/Eclipse Images/EclipseFullText}

%
% Image paths are relative to the compiled document's folder, so every
% document using this file must sit one folder below it.
% ---------------------------------------------------------------------------

% Meeting details, used on covers and in each document's ID block.
\newcommand{\meetingdate}{30 September 2026}
\newcommand{\meetingtime}{12:00}
\newcommand{\meetinglocation}{Festsalen, NMBU}
\newcommand{\meetingline}{\textbf{Meeting:} \meetingdate, \meetingtime, \meetinglocation}

% Header: moon logo on the left, full-text logo on the right. \projectname is
% still set because the cover title uses it.
\projectname{Eclipse NMBU}
\projectlogo{../../../Common/Eclipse Images/EclipseMoon}
\orglogo{../../../Common/Eclipse Images/EclipseFullText}

%
% Image paths are relative to the compiled document's folder, so every
% document using this file must sit one folder below it.
% ---------------------------------------------------------------------------

% Meeting details, used on covers and in each document's ID block.
\newcommand{\meetingdate}{30 September 2026}
\newcommand{\meetingtime}{12:00}
\newcommand{\meetinglocation}{Festsalen, NMBU}
\newcommand{\meetingline}{\textbf{Meeting:} \meetingdate, \meetingtime, \meetinglocation}

% Header: moon logo on the left, full-text logo on the right. \projectname is
% still set because the cover title uses it.
\projectname{Eclipse NMBU}
\projectlogo{../../../Common/Eclipse Images/EclipseMoon}
\orglogo{../../../Common/Eclipse Images/EclipseFullText}

\usepackage{enumitem}

\reporttitle{Plain-Language Explanation}

% ---------------------------------------------------------------------------
% Cover / title page for rapport-class documents.
%
% Usage (after \documentclass{rapport} and before \begin{document}):
%   % ---------------------------------------------------------------------------
% Cover / title page for rapport-class documents.
%
% Usage (after \documentclass{rapport} and before \begin{document}):
%   % ---------------------------------------------------------------------------
% Cover / title page for rapport-class documents.
%
% Usage (after \documentclass{rapport} and before \begin{document}):
%   \input{../cover/rapport-cover}
%   \missionpatch{path/to/patch}      % optional circular mission/team patch;
%                                     % without it, the org logo is shown large
%                                     % at the top instead
%   \coverkind{Vedtekter}             % optional small line above the name
%   \reportsubtitle{Optional subtitle line}
%   \orgname{Eclipse NMBU}{Norwegian University of Life Sciences, Norway}
%   \reportlocation{Ås, Norge}
%   \date{3. oktober 2022}            % optional, defaults to \today
%   \contactinfo{Contact: post@eclipsenmbu.no}  % optional, shown at the bottom
%
% Then, as the first thing inside \begin{document}:
%   \makecoverpage
%
% The large title line shows \projectname (falling back to \reporttitle if
% \coverkind is not set), so a document typically sets both:
%   \projectname{Eclipse NMBU}
%   \coverkind{Vedtekter}
% ---------------------------------------------------------------------------

\makeatletter
\def\@missionpatch{}
\newcommand{\missionpatch}[1]{\def\@missionpatch{#1}}

\def\@coverkind{}
\newcommand{\coverkind}[1]{\def\@coverkind{#1}}

\def\@reportsubtitle{}
\newcommand{\reportsubtitle}[1]{\def\@reportsubtitle{#1}}

\def\@orgnameone{}
\def\@orgnametwo{}
\newcommand{\orgname}[2]{\def\@orgnameone{#1}\def\@orgnametwo{#2}}

\def\@reportlocation{}
\newcommand{\reportlocation}[1]{\def\@reportlocation{#1}}

\def\@contactinfo{}
\newcommand{\contactinfo}[1]{\def\@contactinfo{#1}}

\newcommand{\makecoverpage}{%
  \begingroup
  \thispagestyle{empty}
  \begin{center}
  \ifx\@missionpatch\@empty
    \includegraphics[width=10cm]{\@orglogo}\par
  \else
    \includegraphics[width=6.5cm]{\@missionpatch}\par
  \fi
  \vspace{3.4cm}
  \ifx\@coverkind\@empty
    {\sffamily\Huge\bfseries \@reporttitle \par}
  \else
    {\sffamily\Huge\bfseries \@coverkind \par}
    {\sffamily\Huge\bfseries \@projectname \par}
  \fi
  \ifx\@reportsubtitle\@empty\else
    \vspace{1.6cm}
    {\sffamily\bfseries \@reportsubtitle \par}
  \fi
  \vspace{2cm}
  \ifx\@missionpatch\@empty\else
    \includegraphics[height=1.3cm]{\@orglogo}\par
    \vspace{0.6cm}
  \fi
  \@orgnameone \\
  \@orgnametwo \par
  \vspace{0.6cm}
  \ifx\@reportlocation\@empty
    \@date
  \else
    \@reportlocation, \@date
  \fi
  \par
  \vfill
  \ifx\@contactinfo\@empty\else
    {\small \@contactinfo \par}
  \fi
  \end{center}
  \endgroup
  \newpage
  \pagenumbering{roman}
  \pagestyle{plain}
}
\makeatother

%   \missionpatch{path/to/patch}      % optional circular mission/team patch;
%                                     % without it, the org logo is shown large
%                                     % at the top instead
%   \coverkind{Vedtekter}             % optional small line above the name
%   \reportsubtitle{Optional subtitle line}
%   \orgname{Eclipse NMBU}{Norwegian University of Life Sciences, Norway}
%   \reportlocation{Ås, Norge}
%   \date{3. oktober 2022}            % optional, defaults to \today
%   \contactinfo{Contact: post@eclipsenmbu.no}  % optional, shown at the bottom
%
% Then, as the first thing inside \begin{document}:
%   \makecoverpage
%
% The large title line shows \projectname (falling back to \reporttitle if
% \coverkind is not set), so a document typically sets both:
%   \projectname{Eclipse NMBU}
%   \coverkind{Vedtekter}
% ---------------------------------------------------------------------------

\makeatletter
\def\@missionpatch{}
\newcommand{\missionpatch}[1]{\def\@missionpatch{#1}}

\def\@coverkind{}
\newcommand{\coverkind}[1]{\def\@coverkind{#1}}

\def\@reportsubtitle{}
\newcommand{\reportsubtitle}[1]{\def\@reportsubtitle{#1}}

\def\@orgnameone{}
\def\@orgnametwo{}
\newcommand{\orgname}[2]{\def\@orgnameone{#1}\def\@orgnametwo{#2}}

\def\@reportlocation{}
\newcommand{\reportlocation}[1]{\def\@reportlocation{#1}}

\def\@contactinfo{}
\newcommand{\contactinfo}[1]{\def\@contactinfo{#1}}

\newcommand{\makecoverpage}{%
  \begingroup
  \thispagestyle{empty}
  \begin{center}
  \ifx\@missionpatch\@empty
    \includegraphics[width=10cm]{\@orglogo}\par
  \else
    \includegraphics[width=6.5cm]{\@missionpatch}\par
  \fi
  \vspace{3.4cm}
  \ifx\@coverkind\@empty
    {\sffamily\Huge\bfseries \@reporttitle \par}
  \else
    {\sffamily\Huge\bfseries \@coverkind \par}
    {\sffamily\Huge\bfseries \@projectname \par}
  \fi
  \ifx\@reportsubtitle\@empty\else
    \vspace{1.6cm}
    {\sffamily\bfseries \@reportsubtitle \par}
  \fi
  \vspace{2cm}
  \ifx\@missionpatch\@empty\else
    \includegraphics[height=1.3cm]{\@orglogo}\par
    \vspace{0.6cm}
  \fi
  \@orgnameone \\
  \@orgnametwo \par
  \vspace{0.6cm}
  \ifx\@reportlocation\@empty
    \@date
  \else
    \@reportlocation, \@date
  \fi
  \par
  \vfill
  \ifx\@contactinfo\@empty\else
    {\small \@contactinfo \par}
  \fi
  \end{center}
  \endgroup
  \newpage
  \pagenumbering{roman}
  \pagestyle{plain}
}
\makeatother

%   \missionpatch{path/to/patch}      % optional circular mission/team patch;
%                                     % without it, the org logo is shown large
%                                     % at the top instead
%   \coverkind{Vedtekter}             % optional small line above the name
%   \reportsubtitle{Optional subtitle line}
%   \orgname{Eclipse NMBU}{Norwegian University of Life Sciences, Norway}
%   \reportlocation{Ås, Norge}
%   \date{3. oktober 2022}            % optional, defaults to \today
%   \contactinfo{Contact: post@eclipsenmbu.no}  % optional, shown at the bottom
%
% Then, as the first thing inside \begin{document}:
%   \makecoverpage
%
% The large title line shows \projectname (falling back to \reporttitle if
% \coverkind is not set), so a document typically sets both:
%   \projectname{Eclipse NMBU}
%   \coverkind{Vedtekter}
% ---------------------------------------------------------------------------

\makeatletter
\def\@missionpatch{}
\newcommand{\missionpatch}[1]{\def\@missionpatch{#1}}

\def\@coverkind{}
\newcommand{\coverkind}[1]{\def\@coverkind{#1}}

\def\@reportsubtitle{}
\newcommand{\reportsubtitle}[1]{\def\@reportsubtitle{#1}}

\def\@orgnameone{}
\def\@orgnametwo{}
\newcommand{\orgname}[2]{\def\@orgnameone{#1}\def\@orgnametwo{#2}}

\def\@reportlocation{}
\newcommand{\reportlocation}[1]{\def\@reportlocation{#1}}

\def\@contactinfo{}
\newcommand{\contactinfo}[1]{\def\@contactinfo{#1}}

\newcommand{\makecoverpage}{%
  \begingroup
  \thispagestyle{empty}
  \begin{center}
  \ifx\@missionpatch\@empty
    \includegraphics[width=10cm]{\@orglogo}\par
  \else
    \includegraphics[width=6.5cm]{\@missionpatch}\par
  \fi
  \vspace{3.4cm}
  \ifx\@coverkind\@empty
    {\sffamily\Huge\bfseries \@reporttitle \par}
  \else
    {\sffamily\Huge\bfseries \@coverkind \par}
    {\sffamily\Huge\bfseries \@projectname \par}
  \fi
  \ifx\@reportsubtitle\@empty\else
    \vspace{1.6cm}
    {\sffamily\bfseries \@reportsubtitle \par}
  \fi
  \vspace{2cm}
  \ifx\@missionpatch\@empty\else
    \includegraphics[height=1.3cm]{\@orglogo}\par
    \vspace{0.6cm}
  \fi
  \@orgnameone \\
  \@orgnametwo \par
  \vspace{0.6cm}
  \ifx\@reportlocation\@empty
    \@date
  \else
    \@reportlocation, \@date
  \fi
  \par
  \vfill
  \ifx\@contactinfo\@empty\else
    {\small \@contactinfo \par}
  \fi
  \end{center}
  \endgroup
  \newpage
  \pagenumbering{roman}
  \pagestyle{plain}
}
\makeatother

\coverkind{Bylaw Amendment Proposal (vedtektsendringsforslag)}
\reportsubtitle{Plain-language explanation of the changes: ECL-H2026-CASE-07B}
\orgname{Eclipse NMBU}{Norwegian University of Life Sciences, Norway}
\reportlocation{\meetinglocation, Ås}
\date{\meetingdate, \meetingtime}
\contactinfo{Contact: post@eclipsenmbu.no}

\begin{document}

\makecoverpage
\tableofcontents
\reportmainmatter

\noindent\textbf{Document ID:} ECL-H2026-CASE-07B\\
\meetingline\\
\textbf{Companion to:} the full proposed New Bylaws text, ECL-H2026-CASE-07A, and the voting procedure, ECL-H2026-CASE-07C. This document has no independent legal effect: if it conflicts with ECL-H2026-CASE-07A, the full text controls.

\vsection{Why a restructuring, not a patch}
The Old Bylaws (13 sections) have not kept pace with Eclipse NMBU's growth: the C-suite, membership system, project process and delegation of authority have all grown informally beyond what the Old Bylaws describe. Rather than patch individual sections, the Board proposes a full restructuring into 18 sections across five parts, so that the written Bylaws match how the Association actually operates.

\vsection{What changes}

\textbf{1. C-suite structure.} The Board goes from CEO/CTO/CFO/CMO to CEO/CTO/CBO/CPO.
\begin{itemize}
  \item CFO (finance) and CMO (marketing/external communication) are merged into a single \textbf{CBO: Chief Business Officer}, who holds full, undelegated authority over the Association's economy.
  \item A new \textbf{CPO: Chief People Officer} takes over membership administration, recruitment and internal culture -- work that previously had no dedicated owner.
  \item The \textbf{CTO} becomes the Association's formal \textbf{Deputy Chair (nestleder)}, automatically stepping into the CEO's responsibilities when the CEO is unavailable. This did not exist under the Old Bylaws.
\end{itemize}

\textbf{2. Six membership tiers, not four.} \textbf{Mentor} and \textbf{Alumni} are added as formal tiers alongside Board Member, Core Member, Active Member and Inactive Member. Alumni status is now automatic when a Membership Contract lapses without renewal, rather than undefined.

\textbf{3. Election and Control Committee.} A new three-member committee, elected each autumn, takes over running elections and advance nominations, and monitors Bylaw compliance: work currently done informally by the Board itself. The committee is first elected at a Supplementary General Assembly after the New Bylaws take effect. This General Assembly's own election uses the Old Bylaws' election mechanics (see Case 8).

\textbf{4. Two General Assemblies a year, with different jobs.} The New Bylaws formally distinguish a spring \textbf{Elective General Assembly} (full Board elections) from an autumn \textbf{Supplementary General Assembly} (fills vacancies, elects the Election and Control Committee). The Old Bylaws did not distinguish these: every General Assembly did full elections. This General Assembly is held under the Old Bylaws, so the distinction applies only to future General Assemblies.

\textbf{5. Delegation and signature, formalised.} Under the Old Bylaws, the whole Board jointly held Signature Right. Under the New Bylaws, the \textbf{CEO alone} holds Signature Right and the Association's overall Mandate, but is \textbf{required} to delegate authority to the other Board members in their areas (subject to General Assembly approval), and may delegate Signature Right to them directly. The key idea, carried through this whole document set, is that \textbf{signing on the Association's behalf is itself a responsibility}, not just an administrative permission: see the Responsibility \& Signature Contract, ECL-H2026-RESP-01A/01B.

\textbf{6. Project gates.} Starting a project now formally requires an MCR (Mission Concept Review) and an SRR (System Requirements Review) before the Board approves it, following a modified NASA Project Life Cycle. The Bylaws fix MCR and SRR as mandatory gates; the detailed phase model and any further gate reviews are set out in the Association's governance documents (\S16.3).

\textbf{7. English as operating language.} Eclipse NMBU's operating language becomes English; Norwegian-language documents must be translated to English (and vice versa, on request).

\textbf{8. Privacy.} Members gain an explicit, written right to request deletion of their personal data at any time, with the Association obliged to tell members about this right when they join and when they leave.

\vsection{What does not change}
The Association's name, purpose, independence, self-owning legal status, core ethics/environment principles, whistleblowing process (now clarified to route through the CPO), and the prohibition on proxy voting all carry over unchanged in substance.

\vsection{What this vote actually decides}
Adopting the New Bylaws is a single constitutional decision requiring 2/3 of \emph{all} Eclipse NMBU members (not just those present) to vote in favour: see ECL-H2026-CASE-07C for exactly how that vote is structured alongside any submitted Amendments to this text.

\end{document}
